% Options for packages loaded elsewhere
% Options for packages loaded elsewhere
\PassOptionsToPackage{unicode}{hyperref}
\PassOptionsToPackage{hyphens}{url}
\PassOptionsToPackage{dvipsnames,svgnames,x11names}{xcolor}
%
\documentclass[
  12pt,
  letterpaper,
  man]{apa7}
\usepackage{xcolor}
\usepackage{amsmath,amssymb}
\setcounter{secnumdepth}{-\maxdimen} % remove section numbering
\usepackage{iftex}
\ifPDFTeX
  \usepackage[T1]{fontenc}
  \usepackage[utf8]{inputenc}
  \usepackage{textcomp} % provide euro and other symbols
\else % if luatex or xetex
  \usepackage{unicode-math} % this also loads fontspec
  \defaultfontfeatures{Scale=MatchLowercase}
  \defaultfontfeatures[\rmfamily]{Ligatures=TeX,Scale=1}
\fi
\usepackage{lmodern}
\ifPDFTeX\else
  % xetex/luatex font selection
\fi
% Use upquote if available, for straight quotes in verbatim environments
\IfFileExists{upquote.sty}{\usepackage{upquote}}{}
\IfFileExists{microtype.sty}{% use microtype if available
  \usepackage[]{microtype}
  \UseMicrotypeSet[protrusion]{basicmath} % disable protrusion for tt fonts
}{}
\usepackage{setspace}
\makeatletter
\@ifundefined{KOMAClassName}{% if non-KOMA class
  \IfFileExists{parskip.sty}{%
    \usepackage{parskip}
  }{% else
    \setlength{\parindent}{0pt}
    \setlength{\parskip}{6pt plus 2pt minus 1pt}}
}{% if KOMA class
  \KOMAoptions{parskip=half}}
\makeatother
% Make \paragraph and \subparagraph free-standing
\makeatletter
\ifx\paragraph\undefined\else
  \let\oldparagraph\paragraph
  \renewcommand{\paragraph}{
    \@ifstar
      \xxxParagraphStar
      \xxxParagraphNoStar
  }
  \newcommand{\xxxParagraphStar}[1]{\oldparagraph*{#1}\mbox{}}
  \newcommand{\xxxParagraphNoStar}[1]{\oldparagraph{#1}\mbox{}}
\fi
\ifx\subparagraph\undefined\else
  \let\oldsubparagraph\subparagraph
  \renewcommand{\subparagraph}{
    \@ifstar
      \xxxSubParagraphStar
      \xxxSubParagraphNoStar
  }
  \newcommand{\xxxSubParagraphStar}[1]{\oldsubparagraph*{#1}\mbox{}}
  \newcommand{\xxxSubParagraphNoStar}[1]{\oldsubparagraph{#1}\mbox{}}
\fi
\makeatother


\usepackage{longtable,booktabs,array}
\newcounter{none} % for unnumbered tables
\usepackage{calc} % for calculating minipage widths
% Correct order of tables after \paragraph or \subparagraph
\usepackage{etoolbox}
\makeatletter
\patchcmd\longtable{\par}{\if@noskipsec\mbox{}\fi\par}{}{}
\makeatother
% Allow footnotes in longtable head/foot
\IfFileExists{footnotehyper.sty}{\usepackage{footnotehyper}}{\usepackage{footnote}}
\makesavenoteenv{longtable}
\usepackage{graphicx}
\makeatletter
\newsavebox\pandoc@box
\newcommand*\pandocbounded[1]{% scales image to fit in text height/width
  \sbox\pandoc@box{#1}%
  \Gscale@div\@tempa{\textheight}{\dimexpr\ht\pandoc@box+\dp\pandoc@box\relax}%
  \Gscale@div\@tempb{\linewidth}{\wd\pandoc@box}%
  \ifdim\@tempb\p@<\@tempa\p@\let\@tempa\@tempb\fi% select the smaller of both
  \ifdim\@tempa\p@<\p@\scalebox{\@tempa}{\usebox\pandoc@box}%
  \else\usebox{\pandoc@box}%
  \fi%
}
% Set default figure placement to htbp
\def\fps@figure{htbp}
\makeatother


% definitions for citeproc citations
\NewDocumentCommand\citeproctext{}{}
\NewDocumentCommand\citeproc{mm}{%
  \begingroup\def\citeproctext{#2}\cite{#1}\endgroup}
\makeatletter
 % allow citations to break across lines
 \let\@cite@ofmt\@firstofone
 % avoid brackets around text for \cite:
 \def\@biblabel#1{}
 \def\@cite#1#2{{#1\if@tempswa , #2\fi}}
\makeatother
\newlength{\cslhangindent}
\setlength{\cslhangindent}{1.5em}
\newlength{\csllabelwidth}
\setlength{\csllabelwidth}{3em}
\newenvironment{CSLReferences}[2] % #1 hanging-indent, #2 entry-spacing
 {\begin{list}{}{%
  \setlength{\itemindent}{0pt}
  \setlength{\leftmargin}{0pt}
  \setlength{\parsep}{0pt}
  % turn on hanging indent if param 1 is 1
  \ifodd #1
   \setlength{\leftmargin}{\cslhangindent}
   \setlength{\itemindent}{-1\cslhangindent}
  \fi
  % set entry spacing
  \setlength{\itemsep}{#2\baselineskip}}}
 {\end{list}}
\usepackage{calc}
\newcommand{\CSLBlock}[1]{\hfill\break\parbox[t]{\linewidth}{\strut\ignorespaces#1\strut}}
\newcommand{\CSLLeftMargin}[1]{\parbox[t]{\csllabelwidth}{\strut#1\strut}}
\newcommand{\CSLRightInline}[1]{\parbox[t]{\linewidth - \csllabelwidth}{\strut#1\strut}}
\newcommand{\CSLIndent}[1]{\hspace{\cslhangindent}#1}



\setlength{\emergencystretch}{3em} % prevent overfull lines

\providecommand{\tightlist}{%
  \setlength{\itemsep}{0pt}\setlength{\parskip}{0pt}}



 


\usepackage{csquotes}
\makeatletter
\@ifpackageloaded{caption}{}{\usepackage{caption}}
\AtBeginDocument{%
\ifdefined\contentsname
  \renewcommand*\contentsname{Table of contents}
\else
  \newcommand\contentsname{Table of contents}
\fi
\ifdefined\listfigurename
  \renewcommand*\listfigurename{List of Figures}
\else
  \newcommand\listfigurename{List of Figures}
\fi
\ifdefined\listtablename
  \renewcommand*\listtablename{List of Tables}
\else
  \newcommand\listtablename{List of Tables}
\fi
\ifdefined\figurename
  \renewcommand*\figurename{Figure}
\else
  \newcommand\figurename{Figure}
\fi
\ifdefined\tablename
  \renewcommand*\tablename{Table}
\else
  \newcommand\tablename{Table}
\fi
}
\@ifpackageloaded{float}{}{\usepackage{float}}
\floatstyle{ruled}
\@ifundefined{c@chapter}{\newfloat{codelisting}{h}{lop}}{\newfloat{codelisting}{h}{lop}[chapter]}
\floatname{codelisting}{Listing}
\newcommand*\listoflistings{\listof{codelisting}{List of Listings}}
\makeatother
\makeatletter
\makeatother
\makeatletter
\@ifpackageloaded{caption}{}{\usepackage{caption}}
\@ifpackageloaded{subcaption}{}{\usepackage{subcaption}}
\makeatother
\usepackage{bookmark}
\IfFileExists{xurl.sty}{\usepackage{xurl}}{} % add URL line breaks if available
\urlstyle{same}
\hypersetup{
  pdftitle={ReMis {[}A fun title could go here{]}},
  pdfauthor={Leonie Hagitte},
  colorlinks=true,
  linkcolor={blue},
  filecolor={Maroon},
  citecolor={Blue},
  urlcolor={Blue},
  pdfcreator={LaTeX via pandoc}}


\title{ReMis {[}A fun title could go here{]}}
\author{Leonie Hagitte}
\date{}
\begin{document}
\maketitle


\setstretch{2}
\section{Introduction}\label{introduction}

Structural equation models (SEMs) are widely used in applied research to
test hypotheses about latent constructs and their relationships. In
practice, replication is typically assessed by fitting the same model to
a new dataset and evaluating whether global fit indices remain
acceptable or whether selected parameters retain statistical
significance. Empirically, these approaches often yield ambiguous or
unstable outcomes. Models may exhibit acceptable global fit across
datasets while differing substantially in parameter estimates,
reflecting the non-uniqueness of the model-implied covariance structure.
Conversely, parameters that were statistically significant in an
original study may fail to reach significance in a replication due to
differences in sample size, measurement quality, or sampling
variability. In addition, estimation itself may be unstable. A model
that converges in one dataset may fail to converge in another sample
drawn from the same population, or even in resampled versions of the
original data. Such behavior indicates sensitivity of the model to
sampling fluctuations or weak identification, and implies that
replication may fail at the level of model estimation rather than
parameter comparison. These issues suggest that replication in SEM
requires an explicit framework for distinguishing model fit, parameter
stability, and theoretical stability. We argue that replication of SEMs
can be assessed using the logic of invariance testing: measurement
invariance establishes whether parameters can be meaningfully compared
across datasets, parameter stability provides statistical evidence for
replication, and theoretical stability determines whether the
substantive claim encoded by the model is preserved.

\subsection{Replication}\label{replication}

Replication is the ability to obtain consistent findings when a study is
repeated with new independent data being collected under the same or
sufficiently similar methodological and theoretical conditions as the
original study. In the context of empirical research, a finding can be
understood at different analytical levels: as an empirical regularity,
as an estimand-based claim, or as a theoretical claim about an
underlying mechanism that generates observed patterns. The type or
extent of replication relevant to a finding depends on which of these
levels the finding is defined at.

Radder (1992) characterizes replication as comprising three distinct
forms: (a) material replication; (b) reproducibility of the theoretical
interpretation; and (c) replication of the experimental result through
an alternative material procedure. Schmidt (2009) builds on this
definition and establishes the concepts of direct replication as a
restrictive and narrow form of replication and conceptual replication as
a wider notion of replication.

Direct replication refers to a repetition of the original experimental
or observational procedure aimed at re-estimating the same estimand
under equivalent conditions. Formally, direct replication concerns
approximate agreement in the estimation of the same substantive
quantity:

\[\theta^{*(1)} \approx \theta^{*(2)}.\]

Under finite sampling, exact numerical equality of parameter estimates
is generally implausible. Thus, direct replication refers to a
repetition of the original experimental or observational procedure aimed
at re-estimating the same estimand under equivalent conditions. This in
turn relates to non-experimental methods as when a new study estimates
the same estimand under conditions that are equal to the original,
operationalizing the variables in the same way, with differences limited
to incidental sources of variation such as sample, time, or setting.
Under this definition, deviations from the original estimate should be
small, but exact numerical equality is not required.

By contrast, conceptual replication concerns the stability of a
theoretical claim under changes in operationalization. This relates to
Radder (1992), his b and c compartments of a replication, so that the
same theoretical claim or construct is supported using a different
estimand that is theoretically analogous to the original. This involves
changes in operationalization, such as different independent or
dependent measures, while retaining the underlying conceptual target.
Conceptual replication can then be expressed through a mapping \(\phi\)
such that

\[C(\theta^{(1)}) = C(\phi(\theta^{(2)})), \]

where \(C(\cdot)\) denotes the theoretical claim supported by the
estimand. In this case, replication is evaluated at the level of
theoretical equivalence rather than parameter identity.

\subsection{Replicating Structural Equation
Models}\label{replicating-structural-equation-models}

According to Kaplan (2009), SEM comprises a class of methodologies
designed to express a hypothesized model through summary statistics
derived from empirical observations, parameterized by a reduced set of
structural parameters. As SEMs are also considered to be a way to
formulate social science theories(De Leeuw and Berk 2009), for a SEM to
be deemed replicated, the encompassed theory needs to hold in the data,
either globally or regarding certain parameters.

For models involving continuous latent variables, the associated
hypotheses typically pertain to the means, variances, and covariances of
the observed variables. Hypotheses may also concern effect sizes,
ordering of parameter magnitudes, or specific parameter values. Because
SEMs often contain many estimated parameters, replication should
generally focus on those outcomes that are substantively relevant to the
hypothesis under study.

\subsubsection{Global model fit versus parameter-level
replication}\label{global-model-fit-versus-parameter-level-replication}

Global model fit concerns the adequacy of the model-implied covariance
structure. Let \(\boldsymbol{\Sigma}(\theta)\) denote the covariance
matrix implied by the model parameters \(\theta\), and let
\(\mathbf{S}^{(d)}\) denote the observed covariance matrix in dataset
\(d\). Global fit evaluates whether

\[\boldsymbol{\Sigma}(\theta) \approx \mathbf{S}.\]

However, good global fit in both datasets does not imply replication at
the level of parameters. Multiple parameter configurations may produce
similar model-implied covariance structures. Global fit therefore
concerns the adequacy of the model-implied distribution, whereas
replication at the level of parameters concerns the stability of
theoretically relevant estimands.

Accordingly, replication in SEM should be evaluated with respect to
those parameters that directly encode the substantive claim under
investigation. Let \(\theta^*\) denote the subset of theory-relevant
model parameters. Replication at the parameter level may then be
expressed as

\[
\theta^{*(1)}
\approx
\theta^{*(2)}.
\]

This formulation captures stability of theory-relevant estimands rather
than merely adequacy of the model-implied covariance structure. At the
same time, parameter stability should not be equated with replication
success itself. Replication ultimately concerns the stability of the
substantive conclusions supported by these parameters.

If replication is viewed as a comparison of theory-relevant parameters
across independent datasets, both datasets may be treated as
realizations of the same underlying population model. The resulting
problem is therefore one of systematic parameter comparison across
datasets. In the following sections, we argue that the logic of
measurement and structural invariance provides a natural methodological
framework for addressing this problem.

\subsection{Invariance Testing}\label{invariance-testing}

Measurement invariance (MI) refers to the statistical property that a
measurement instrument assesses the same latent construct equivalently
across different groups. In other words, the measurement model operates
in the same manner across groups. When MI is established, meaningful
comparisons of the underlying construct between groups are justified.

From a statistical modelling perspective, MI refers to invariance of the
measurement parameters across groups.
\[ \theta_M^{(g)} = \theta_M \quad \forall g ,\] where \(\theta_M\)
denotes the set of measurement parameters, including factor loadings,
intercepts, and residual variances.

Structural invariance (SI), by contrast, refers to equality of
structural coefficients across groups, that is, to the requirement that
the same association between latent factors exists across groups.
\[\theta_S^{(g)} = \theta_S \quad \forall g,\] where \(\theta_S\)
denotes the set of structural parameters governing the relationships
among latent variables.

Traditionally, MI and SI are characterized via increasingly strict
levels of invariance that each add further restrictions on statistical
parameters across groups. In the classical view, MI is attained only
when configural, metric, scalar, and residual invariance hold. These
correspond to successive equality restrictions on factor loadings,
intercepts, and residual variances across groups. More recent
perspectives conceptualize invariance more generally as the absence of
systematic moderation of model parameters by variables external to the
model. Formally, if model parameters depend on group membership through

\[ \theta^{(g)} = \theta + \Delta^{(g)}, \]

then invariance corresponds to the condition

\[ \Delta^{(g)} = 0 \quad \forall g. \]

Covariance invariance analyses evaluate whether covariances among latent
factors are equivalent across groups, whereas SI analyses examine
whether the structural relationships among latent factors are equal
across groups. Together, these analyses assess whether relationships
among latent variables remain stable across groups or vary
systematically as a function of group membership.

\subsection{Using Invariance Testing to replicate
SEMs}\label{using-invariance-testing-to-replicate-sems}

The central argument of this paper is that the logic of invariance
testing can be repurposed as a methodological framework for assessing
replication in SEM. Measurement invariance is not itself a theory of
replication. Rather, it provides a statistical procedure for evaluating
whether theory-relevant parameters remain stable under a specified
source of variation. In traditional applications, this variation is
induced by group membership. In replication research, the relevant
source of variation is independent data collection. The methodological
role of invariance testing therefore remains unchanged, while the
substantive interpretation shifts from group comparison to replication
assessment.

Whereas measurement invariance is conditional on group definitions,
replication is conditional on independent datasets. Nevertheless, both
involve the assessment of parameter stability under variation. This
makes the relation between replication and invariance testing explicit:
both evaluate whether theory-relevant parameters remain stable when
conditions change, differing primarily in the source of that
variation.\\
Replication of SEMs across an original and a replication sample should
therefore proceed within a sequential multiple-group manner.

Replication, as defined above, refers to the stability of
theory-relevant parameters across independent datasets and is not
conditional on measurement invariance. However, in the context of SEM,
the empirical assessment of structural parameter stability requires that
the underlying measurement model operates equivalently across datasets.
Without such equivalence, observed differences in structural parameters
may reflect differences in measurement rather than differences in the
underlying relations of interest. Measurement invariance is therefore a
necessary condition for the valid interpretation of structural
replication.

The analogy to invariance testing applies primarily to the statistical
assessment of parameter stability. As in measurement invariance
research, evidence against exact equality does not necessarily imply
substantively meaningful non-equivalence. Accordingly, the outcome of
invariance tests should be interpreted with respect to the theoretical
claims represented by the constrained parameters. This suggests a
distinction between the statistical assessment of replication and its
substantive interpretation. Measurement invariance provides the
methodological framework, parameter stability constitutes the
statistical target, and theoretical stability serves as the criterion
for replication. The former concerns the stability of theory-relevant
parameters across datasets, whereas the latter concerns whether any
observed differences alter the theoretical conclusions supported by the
model.

First, measurement invariance between the two samples has to be
established:

\[\theta_M^{(1)} = \theta_M^{(2)}. \]

At a minimum, sufficient measurement invariance must be established for
the structural parameters of interest to be meaningfully compared across
datasets. Conditional on sufficient measurement invariance, structural
replication is assessed by comparing an unconstrained multiple-group
model, in which the structural parameters of interest are freely
estimated across datasets, with a constrained model in which they are
set to be equal:

\[ H_0: \theta_S^{*(1)} = \theta_S^{*(2)}. \]

This comparison could be evaluated with a likelihood-ratio statistic. If
a model in which the parameters are constrained to be equal across
groups fits the data no worse than a model in which they are freely
estimated, this is consistent with structural stability of those
parameters. Conversely, if the constrained model fits the data worse,
this indicates that the parameters differ across groups. At the same
time, such tests should be interpreted cautiously because they are
sensitive to sample size and do not by themselves capture uncertainty in
a substantively calibrated way.

\subsection{Quantifying and Evaluating
Discrepancies}\label{quantifying-and-evaluating-discrepancies}

Evidence against exact equality does not necessarily imply replication
failure. As in the measurement invariance literature, statistically
detectable differences may have negligible substantive consequences
{[}NYE Reference{]}. The question therefore shifts from whether
parameters differ to whether the observed differences alter the
theoretical claim supported by the model.

This perspective is consistent with the earlier formulation of
conceptual replication, where replication was evaluated through
preservation of a theoretical claim rather than equality of estimands.
The same logic can be extended to direct replication. Let

\[
C(\theta)
\]

denote the theoretical claim implied by a parameter configuration
\(\theta\). Direct replication may then be interpreted as preservation
of the same theoretical claim across independent datasets:

\[
C\left(\theta^{*(1)}\right) = C\left(\theta^{*(2)}\right).
\]

Under this formulation, parameter stability constitutes evidence for
replication, but replication itself is evaluated at the level of
theoretical conclusions. Consequently, statistically detectable
parameter differences need not imply replication failure if the
substantive claim represented by those parameters remains unchanged.

This formulation unifies direct and conceptual replication within a
common framework. For conceptual replication, preservation of a
theoretical claim is evaluated across theoretically related estimands:

\[
C\left(\theta^{(1)}\right) = C\left(\phi\left(\theta^{(2)}\right)\right),
\]

where \(\phi\) denotes a mapping between different operationalizations.
By contrast, direct replication concerns preservation of the same
theoretical claim for the same estimand across independent datasets:

\[
C\left(\theta^{*(1)}\right) = C\left(\theta^{*(2)}\right).
\]

This implies that replication of SEMs is best understood as a graded,
parameter- and claim-relative process.\\
Because theoretical claims are evaluated through model parameters,
assessment of theoretical stability requires a corresponding assessment
of parameter stability. Let

\[ \theta^* = \{\theta_1^*, \dots, \theta_r^*\} \]

denote the set of parameters relevant to the theoretical claim of
interest. Parameter stability can then be represented as the vector

\[ \hat{S} = \left( \mathbb{I}\left(\left|\hat\theta_j^{*(1)} - \hat\theta_j^{*(2)}\right| \le \epsilon_j\right) \right)_{j=1}^r, \]

where \(\mathbb{I}(\cdot)\) is an indicator function and \(\epsilon_j\)
is a substantively justified tolerance for parameter \(j\). A parameter
is said to be empirically \emph{stable} if

\[\left|\hat\theta_j^{*(1)} - \hat\theta_j^{*(2)}\right| \le \epsilon_j.\]

Direct and conceptual replication therefore differ in the degree of
operational similarity between studies, but both may be understood as
instances of theoretical stability. In direct replication, stability is
assessed for the same estimand across datasets, whereas in conceptual
replication stability is assessed across theoretically related estimands
linked through \(\phi\).

Full replication corresponds to preservation of all theory-relevant
claims across datasets. Partial replication corresponds to preservation
of only a subset of those claims. Parameter stability, as represented by
the criterion above, provides empirical evidence regarding replication
but should not be interpreted as identical to replication itself. This
perspective renders approximate and partial replication theoretically
grounded rather than ad hoc, because limited deviations in parameter
estimates can be interpreted as meaningful quantitative differences
rather than categorical failures of replication, paralleling
contemporary approaches to approximate invariance.

\subsection{Practical Evaluation of Replication
Outcomes}\label{practical-evaluation-of-replication-outcomes}

The preceding sections established a distinction between parameter
stability and theoretical stability. Measurement invariance provides the
methodological framework for comparing parameters across datasets,
whereas replication is ultimately evaluated with respect to the
theoretical claims supported by those parameters. The remaining question
is therefore how observed parameter discrepancies should be quantified
and interpreted in practice.

Consistent with Nye et al.'s consequence-based perspective on
measurement nonequivalence, replication assessment may focus on the
substantive consequences of parameter discrepancies rather than their
mere statistical presence. In practical applications, replication
assessment therefore proceeds in three stages: (a) establish measurement
comparability, (b) assess parameter stability, and (c) evaluate whether
any remaining discrepancies alter the theoretical claim of interest. The
framework proposed here does not prescribe a unique discrepancy metric.
Rather, the choice of metric should reflect the substantive consequences
of the observed parameter differences. In this sense, replication
assessment may be viewed as consequence-based: the relevant question is
not merely how much parameter estimates differ, but how much the
theoretical implications of the model change as a result.

Several operationalizations are possible. Discrepancies may be expressed
as standardized parameter differences, relative deviations from the
original estimate, changes in latent covariance structures, or changes
in the model-implied covariance matrix. The common principle is that
discrepancy measures should be interpreted with respect to their
consequences for the theoretical claim under investigation.

\subsubsection{Example 1: Full Replication Despite Parameter
Differences}\label{example-1-full-replication-despite-parameter-differences}

Consider a structural path relating a latent predictor to a latent
outcome. In the original study, the standardized path coefficient is

\(\hat{\beta}^{(1)} = .50,\) whereas the replication study yields
\(\hat{\beta}^{(2)} = .45.\) The absolute discrepancy is
\(\left|\hat{\beta}^{(1)}-\hat{\beta}^{(2)}\right|=.05.\)

A relative discrepancy measure may be defined as

\[
W=
\frac{\left|\hat{\beta}^{(1)}-\hat{\beta}^{(2)}\right|}
{\left|\hat{\beta}^{(1)}\right|}.
\] Substituting the observed values gives \(W= \frac{.05}{.50} = .10.\)

Thus, the replication estimate differs by approximately 10\% from the
original estimate. The discrepancy is therefore small relative to the
magnitude of the original effect. Despite this discrepancy, both studies
support the same substantive conclusion: the predictor exhibits a
positive and moderately sized effect on the outcome. Under the framework
proposed here, replication would therefore be considered successful
because the theoretical claim is preserved:

\[
C\left(\theta^{*(1)}\right) = C\left(\theta^{*(2)}\right).
\]

\subsubsection{Example 2: Approximate
Replication}\label{example-2-approximate-replication}

Now consider a smaller original effect \(\hat{\beta}^{(1)}=.10,\) with
the replication estimate \(\hat{\beta}^{(2)}=.05.\) The absolute
discrepancy is again \(.05,\) but the relative discrepancy becomes

\[
W= \frac{.05}{.10} = .50.
\]

Thus, approximately half of the original effect has disappeared. The
absolute discrepancy is identical to that in Example 1, yet its
substantive consequences are considerably larger. This illustrates why
absolute parameter differences alone may be misleading as indicators of
replication success. Depending on the theoretical context, the
underlying claim may still be regarded as preserved, but with weaker
empirical support. This situation illustrates approximate replication:
parameter differences exist, yet the theoretical implication remains
largely intact.

\subsubsection{Example 3: Partial
Replication}\label{example-3-partial-replication}

Suppose a theory predicts three relations among latent variables.

{\def\LTcaptype{none} % do not increment counter
\begin{longtable}[]{@{}lll@{}}
\toprule\noalign{}
Relation & Study 1 & Study 2 \\
\midrule\noalign{}
\endhead
\bottomrule\noalign{}
\endlastfoot
A → B & .70 & .68 \\
A → C & .40 & .41 \\
B → C & .50 & .20 \\
\end{longtable}
}

A discrepancy measure based on the latent relation matrix may be defined
as

\[
W=
\max
\left|
R^{(1)}-R^{(2)}
\right|.
\]

For the example above \(W=.30.\), indicating that at least one
theoretically relevant relation differs substantially across studies.
Consequently, some theoretical claims are preserved while others are
not. Under the proposed framework, this outcome would be interpreted as
partial replication rather than complete success or complete failure.

\subsubsection{Example 4: Replication
Failure}\label{example-4-replication-failure}

Finally, Consider an original study yielding \(\hat{\beta}^{(1)}=.45,\)
whereas the replication study yields \(\hat{\beta}^{(2)}=-.05.\)
Although both estimates may be statistically distinguishable from zero,
the substantive interpretation changes fundamentally. The original study
supports a positive relation between the constructs, whereas the
replication study suggests no meaningful relation or even a weak
negative association.

In this case,
\(C\left(\theta^{*(1)}\right)\neq C\left(\theta^{*(2)}\right),\) because
the theoretical conclusion implied by the model is no longer preserved.
Replication would therefore be considered unsuccessful.

These examples illustrate that replication assessment depends not only
on the magnitude of parameter differences but also on their consequences
for the theoretical claims represented by the model. Consequently,
discrepancy measures should be interpreted with respect to the
substantive conclusions they support rather than in isolation.

\section{Discussion}\label{discussion}

It follows that models in their entirety should not be treated as the
primary objects of replication. Rather, models serve as entry points for
evaluating the stability of the theoretical commitments they encode.
This shifts attention away from global model fit toward the robustness
of substantively meaningful parameters. Because parameter stability
under model misspecification may reflect shared misspecification rather
than stability of the underlying relations, global fit is a necessary
but not sufficient condition for meaningful interpretation of
replication. At the same time, imperfect global fit does not necessarily
invalidate replication at the level of specific theory-relevant
parameters. Measurement quality also becomes central to replication
validity, because instability across datasets may reflect violations of
measurement equivalence rather than instability of the structural
relations themselves. Replication claims that do not explicitly address
measurement equivalence therefore remain ambiguous.

Conceptual replication has a formal structure under this perspective as
it concerns stability of higher-level theoretical claims under
deliberate changes in operationalization. This can be understood as
stability of \(C(\cdot)\) under transformations of the estimand space,
rather than equality of the same parameter across datasets. Replication
success therefore becomes explicitly conditional on the estimand, the
theoretical scope, and the operational context, requiring precise
statements about what exactly has or has not replicated. The distinction
between invariance testing and replication may therefore be understood
as one of scope rather than kind. Invariance evaluates stability across
groups within a dataset, whereas replication evaluates stability across
independent datasets or study instantiations, but both address the same
underlying question of robustness under variation.

\section*{References}\label{references}
\addcontentsline{toc}{section}{References}

\protect\phantomsection\label{refs}
\begin{CSLReferences}{1}{1}
\bibitem[\citeproctext]{ref-DeLeeuw2009}
De Leeuw, J., and R. Berk. 2009. {``{Preface}.''} In \emph{Structural
Equation Modeling: Foundations and Extensions}, 2nd ed., edited by J. De
Leeuw and R. Berk. SAGE Publications, Inc.
\url{https://doi.org/10.4135/9781452226576}.

\bibitem[\citeproctext]{ref-Kaplan2009}
Kaplan, David. 2009. \emph{Structural {Equation} {Modeling} (2nd Ed.):
{Foundations} and {Extensions}}. 2nd ed. Edited by J. De Leeuw and R.
Berk. SAGE Publications, Inc.
\url{https://doi.org/10.4135/9781452226576}.

\bibitem[\citeproctext]{ref-Radder1992}
Radder, Hans. 1992. {``Experimental Reproducibility and the
Experimenters' Regress.''} \emph{PSA: Proceedings of the Biennial
Meeting of the Philosophy of Science Association} 1992: 63--73.
\url{http://www.jstor.org/stable/192744}.

\bibitem[\citeproctext]{ref-Schmidt2009}
Schmidt, Stefan. 2009. {``Shall We {Really} Do It {Again}? {The}
{Powerful} {Concept} of {Replication} Is {Neglected} in the {Social}
{Sciences}.''} \emph{Review of General Psychology} 13 (2): 90--100.
\url{https://doi.org/10.1037/a0015108}.

\end{CSLReferences}




\end{document}
